%% ==========================================================================
%% SECTION 2: BACKGROUND
%% ==========================================================================
\section{Background \& Motivation}\label{sec:background}

A language model maps input tokens to output tokens, producing
internal signals (entropy, attention patterns, log\-probabilities)
as computational byproducts. The training corpus contains both
truth-tracking institutional text and misinformation.
Probing how a model trained on this mixture responds yields four canonical
query types that span the verification space: \emph{adversarial truths} (true
but implausible), \emph{shattered lies} (no training-data grounding),
\emph{deceived lies} (fabrications that pattern-match plausible language), and
\emph{confused truths} (counterintuitive facts).

A model trained on this corpus inherits its mixed regularities: base models
fabricate substantially on factual benchmarks, yet also encode internal
representations of correctness that can be elicited under the right
conditions~\cite{burns2023discovering}. Consequently, the model is neither reliably honest
nor uniformly deceptive.
This uncertainty makes verification necessary:
the challenge is determining which outputs require scrutiny and at what cost.

\subsection{Testimony \& Telemetry}

\fakepara{Architecture primer.}
A transformer generates text one token at a time. At
each step it produces a probability distribution over its vocabulary, selects
a token, and conditions the next step on all prior output. Each step passes
through a stack of layers producing \emph{attention patterns} (weighted
relevance scores between positions) and intermediate representations
to generate a final text output, which we call its \emph{testimony}. The
\emph{entropy} of the output distribution, the attention patterns, and the
\emph{log-probabilities} of selected tokens are byproducts of this
computation, which we call its \emph{telemetry}.

From a verification cost perspective, testimony is cheap to produce and inspect but unreliable;
telemetry is more expensive to export and process but harder to game (under current training
regimes (\autoref{sec:threat})), because it is a byproduct of the computation rather
than a product of the optimization objective.
Existing models expose only its testimony and discard its telemetry.

\subsection{The Text Channel Cannot Self-Verify}
\label{sec:bg_eval}

% Current models cannot
% be trusted to honestly report their own epistemic state through the text channel.
To demonstrate the inadequacy of text-only verification, we use a taxonomy of
query types defined by construction: \emph{knowable} queries have a determinate, publicly verifiable
answer (e.g., national capitals, well-documented scientific facts), whereas
\emph{unknowable} queries have no correct answer for any model, because the
entity does not exist, the event has not happened, or the information is
private. ``Knowable'' is a property of the query, not a claim about a
particular model's weights; a model may still fail a knowable query.
We report each model's self-confidence
score for the self-reported confidence baseline, the cheapest verification strategy
available to a bounded supervisor. In this baseline, the supervisor
asks the model ``How confident are you?'' after generation and tests whether the
model can introspect and report its epistemic state through text.

\begin{figure}[t]%
\centering%
\includegraphics[width=\columnwidth]{figures/exp27_confidence_distributions.pdf}%
\caption{Model self-reported confidence distribution}
\label{fig:confidence_dist}%
\end{figure}

\autoref{fig:confidence_dist} shows that self-reported confidence is high
almost everywhere: models report a mean of 0.98 on knowable queries and 0.73
on unknowable ones, where the correct report is that no answer exists. The
gap orders the two classes better than chance (AUC 0.66--0.92 across
architectures), but the scale is compressed into its top end, so a threshold
that admits most correct answers also admits most fabrications. A user who
asks ``how confident are you?'' is told 100 either way, and an adversary need
only stay inside the band the model already occupies. Under bounded oversight,
a text-only supervisor cannot distinguish uncertainty from fabrication.
